% !PRINTED 528
\weberpage{528}

If, however, we eliminate from the equations (2) and (3) the powers of $\Theta$, there results a relation of the form
\weq{6}{\omega = k_1\omega_1 + k_2\omega_2 + \cdots + k_n\omega_n.}

Every number and every functional of the field $\Omega$ can therefore be represented in the form (6), if one puts rational numbers or functionals for $k_1, k_2, \ldots, k_n$. For a given $\omega$ this representation is possible [because of (5)] in only one way, and $\omega$ is a number when the $k_1, k_2, \ldots, k_n$ are numbers, but a functional when functionals also occur among the $k$.

A system of numbers such as
\[
\omega_1, \omega_2, \ldots, \omega_n,
\]
we call a \gesperrt{basis of the field $\Omega$}. The powers of $\Theta$ also form such a basis:
\[
1, \Theta, \Theta^2, \ldots, \Theta^{n-1}.
\]

If we denote by $\omega_{r,1}, \omega_{r,2}, \ldots, \omega_{r,n}$ the numbers conjugate to $\omega_r$, then the square of the determinant
\weq{7}{\Delta(\omega_1, \omega_2, \ldots, \omega_n) = (\Sigma \pm \omega_{1,1}\omega_{2,2}\ldots\omega_{n,n})^2}
is a symmetric function of the conjugate values $\Theta_1, \Theta_2, \ldots, \Theta_n$ and is consequently a rational number.

This number is called the \gesperrt{discriminant of the system} $\omega_1, \omega_2, \ldots, \omega_n$. In particular,
\weq{8}{\Delta(1, \Theta, \Theta^2, \ldots, \Theta^{n-1}) = \begin{vmatrix} 1, & \Theta_1, & \ldots, & \Theta_1^{\,n-1}\\ 1, & \Theta_2, & \ldots, & \Theta_2^{\,n-1}\\ \multicolumn{4}{c}{\cdots\cdots\cdots\cdots\cdots\cdots}\\ 1, & \Theta_n, & \ldots, & \Theta_n^{\,n-1} \end{vmatrix}^2}
is the discriminant of the equation (1) (Bd.\ I, \S.\,46), for which, if $N$ is the symbol for the norm, one can also put
\weq{9}{(-1)^{\frac{n(n-1)}{2}} N f'(\Theta)}
This number is therefore certainly different from zero. By the multiplication theorem of determinants (Bd.\ I, \S.\,27) there results from (3) and (7):
\weq{10}{\Delta(\omega_1, \omega_2, \ldots, \omega_n) = H^2\Delta(1, \Theta, \Theta^2, \ldots, \Theta^{n-1}),}
from which one concludes that the discriminant of a basis of $\Omega$ is always different from zero. Since $H$ is a rational number, it follows that the ratio of the discriminants of different bases is the square of a rational number, and that the
%CONT>>
